\documentclass[openany,fontset=fandol]{ctexbook}
\usepackage{amsmath,amssymb}
\usepackage[unicode,hidelinks]{hyperref}
\newcommand{\note}[1]{\footnote{译注：#1}}
\begin{document}
\title{代数几何基础第五卷预备笔记\\第12节：维数条件与推广计划}
\author{据用户上传英译编辑稿译出}
\date{2026年11月16日}
\maketitle
\chapter*{范围与校订状态}
本批连续译出底本第46页“Proposition 6.1”起的五段，止于“the statements that we could perhaps give in this connection (a ce proper)”。下一接点是同页“Let us again assume that”段。所译为第12节对第6、7、8节的推广说明与编写计划，不是这些条目的完整重述和证明。本稿独立保存，未并入主稿；既有缺文及待核事项继续保留。
\chapter*{§12 将上述结果推广到线性截面（续）}
\section*{第6节与第7节}
命题6.1中，应读作
\[
\dim f(X_i)\geq m+1,
\]
并在后文读作
\[
\dim f(Z)\leq m-1.
\]
\note{本段只是原作者对6.1中两处条件的修改指示，不是命题6.1的完整陈述；其余假设并未在此取消。}

7.1中的一般论述，对于任意$m$仍可照原样保留。7.2、7.3作必要的相应修改即可（请留意7.2中的记号$m$，它在那里容易造成混淆）。另一方面，在7.4的证明中，不再需要逐步推进，而可以直接取一个余维$m=n$的线性截面。\note{底本在“一般论述”处保留“laius (speech) [Fr or Latin] (speech) [. 69]”，在“作必要的相应修改”处用拉丁语“mutatis mutandis”。这里的$m=n$是该证明中选择截面余维的指示，不能解释成本节所有情形下两个符号恒等。}

\newpage
\section*{第8节的修改与编写取舍}
在8.2中，把其中的超额维数条件改为
\[
\dim Y_\xi>\dim X-m
\qquad\text{〔底本标注：字迹不清，请询问A.G.〕},
\]
并把假设$\dim f(X_i)>0$替换为$\dim f(X_i)\geq m$。8.2后续部分作类似修改。由于8.3给的是一个例子，改动它没有意义，所以保留$m=1$。\note{本段底本把被替换的条件写作$\dim Y_\xi>\dim X$，但前文8.2实际印作$\dim Y_\xi>\dim X-1$，故正文以“其中的超额维数条件”明确指回原条目，原写法留于此。目标条件之后的“illegible, ask AG”仍保留，不将其当作已核清的公式。此外，8.2原假设的阈值本身也留有难辨标记，本段给出的$>0$不能消除那里的原稿疑点。}

我把这件事作为练习留给你：找出对应于8.4、8.5、8.6（第30、31、32页）、适用于任意$m$的恰当陈述。除非你愿意，否则不必做这道练习。\note{英译者在“你”后注“[Dieudonn\'e or Blass] [Tr]”，受话者并未确定。所引页码按底本保留，不能据此认定为本译稿页码。本句是原作者的编写建议，不是向当前读者或译者发出的任务指令。}

我认为，除8.6以外，第8节的讨论基本上都能推广到任意$m$的线性截面情形。把这件事正式写出来，无疑会是一项相当冗长而繁琐的练习。我得承认，我还不知道有哪一种应用在本质上依赖于对这一更一般情形的分析，因此我们并非一定要把这些讨论收入这部《基础》。另一方面，经验表明，在更一般的背景中写作，往往会迫使人更好地把问题拆解开来，从而更清楚地理解整套东西，而且通常并不会增加多少麻烦。此外，在这里这样一个真正的几何背景中，做一些形式表述方面的练习也并无坏处。当然，完全不能排除有朝一日我们会用到或需要它，到时能找到现成的论述，会很高兴。不过，关于这件事的决定仍完全由你作出；我只把我们或许能就此给出的陈述概括一下。\note{本段保留作者试探性的判断和编写取舍。“拆解”对应法语“d\'evisser”，“整套东西”对应“le fourbis”，“通常并不会增加多少麻烦”对应“sans beaucoup plus de mal”。原文的“syntax exercises”在此按编写语境译出，而非补入新的数学结论。上一段把8.6列入可选练习，本段却将其排除于总体推广判断之外；两段的区别按原文保留，未擅自统一。}
\end{document}
